\documentclass[10pt,a4paper]{amsart}
\usepackage[margin=19mm]{geometry}
\usepackage{amsmath,amssymb,mathtools}
\usepackage[scr=rsfs,cal=euler]{mathalfa}
\usepackage{xfrac}
\usepackage[unicode,hidelinks,bookmarks=false]{hyperref}
\newcommand{\ie}{i.e.\ }
\newcommand{\cf}{cf.\ }
\newcommand{\resp}{respectively\ }
\newcommand{\Subsection}{\subsection}
\providecommand{\nobreakdash}{\nobreak\hbox{-}}
\setlength{\emergencystretch}{2em}
\multlinegap0pt
\usepackage{dsfont}
\usepackage{amssymb}
\usepackage[shortlabels]{enumitem}
\usepackage{mathtools}
\usepackage{bm}
\usepackage{tcolorbox}
\newtheorem{defn}{Definition}
\newtheorem{prop}{Proposition}
\def\bbk{\mathbb{K}}
\def\defeq{:=}
\def\fq{\bbk}
\def\fqttt{\fq\!\left(\!\left(t^{-1}\right)\!\right)}
\def\fttt{\fq\!\left(\!\left(t^{-1}\right)\!\right)}
\title{Sums of two continued fractions with growing elements and their finite fields analogues}
\author{Dmitry Gayfulin, Erez Nesharim}
\date{Source: arXiv:2511.11832v2 (CC BY 4.0)}
\begin{document}
\maketitle

\noindent\emph{Source and licence.} Sums of two continued fractions with growing elements and their finite fields analogues. Version arXiv:2511.11832v2 (CC BY 4.0), distributed by arXiv under the Creative Commons Attribution 4.0 International licence, http://creativecommons.org/licenses/by/4.0/, as recorded in the arXiv licence metadata for this exact version and shown on the abstract page https://arxiv.org/abs/2511.11832v2. Full licence notice: \texttt{licenses/arxiv-2511.11832-CC-BY-4.0.txt}.\par\smallskip

\noindent\textit{Editorial scope.} Counted unit: the article's prop beta\_gamma\_def (source0.tex lines 1017--1025). The article's complete original proof of it is delivered with it (source0.tex lines 1026--1094, one \texttt{proof} environment of 69 lines). No mathematics is written, repaired or re-derived: the statement and the proof are the article's own lines, in the article's own notation, and no retained line is substituted or re-flowed.  Retained uncounted as the source-local context the proof needs: lines 335--339; lines 341--343; lines 673--707.  Nothing was omitted from any retained span, so the package carries no omission record.  No retained line contains a citation, so the wrapper carries no bibliography and no bibliography entry.  No retained line contains a figure environment, an image-inclusion command or drawing code, so no image is delivered and the package declares no asset dependency.  Editorial formatting only, in addition: an AMS \texttt{amsart} preamble that restates the article's own declarations and macros used by the retained lines, namely the environments \texttt{defn}, \texttt{prop} on separate counters, together with the standard packages those lines need.  The retained environments print in this document as Definition 1, Proposition 1 in order of appearance, whereas the article prints the same items with the numbering of its own full document.  The complete native source of the article as distributed in the arXiv e-print for this exact version is bundled unchanged as \texttt{sources/189152/source0.tex}.
\begin{equation}\label{eq:common3}
\left|\alpha-\frac{p_n}{q_n} \right|=\frac{1}{|q_n|\,|q_{n+1}|}\,.
%\frac{1}{2^{\deg(q_{n+1})+\deg(q_{n})}}=
%\frac{1}{2^{2\deg(q_n)+\deg(a_{n+1})}}.  
\end{equation}

\begin{equation}\label{eq:common4}
    \left|\alpha-\frac{p}{q} \right|<\frac{1}{|q|^2} \quad\Longrightarrow\quad \frac pq=\frac {p_n}{q_n}\,\textrm{ for some $n\geq0$}\,. 
\end{equation}

\begin{defn}\label{def:shulga_an}
For any $\alpha\in\fqttt$ such that $a_0(\alpha)=0$ define the sequences $\{b_n\}_{n\geq1}$ and $\{c_n\}_{n\geq1}$, and the Laurent series $\beta$ and $\gamma$ as follows:
\begin{enumerate}
\item
For every $n\geq0$, having the integers $b_j$ and $c_j$ defined for every $1\leq j\leq n$, if
\begin{equation*}
\alpha=[0;b_1,\ldots,b_{n}]+[0;c_1,\ldots,c_{n}]\,,
\end{equation*}
then let
$$\beta\defeq[0;b_1,\ldots,b_{n}]\,, \quad \textrm{and}\quad \gamma\defeq[0;c_1,\ldots,c_{n}]\,,$$
and stop. Otherwise, let
\begin{equation*}
b_{n+1} \defeq a_{n+1}(\alpha-[0;c_1,\ldots,c_{n}])\,.
\end{equation*}
If 
\begin{equation*}
\alpha=[0;b_1,\ldots,b_{n+1}]+[0;c_1,\ldots,c_{n}]\,,
\end{equation*}
then let 
$$\beta\defeq[0;b_1,\ldots,b_{n+1}]\,, \quad \textrm{and}\quad \gamma\defeq[0;c_1,\ldots,c_{n}]\,,$$ 
and stop. Otherwise, let 
\begin{equation*}
c_{n+1} \defeq a_{n+1}(\alpha-[0;b_1,\ldots,b_{n+1}])\,.
\end{equation*}
%If
%\begin{equation*}
%\alpha=[0;b_1,\ldots,b_n]+[0;c_1,\ldots,c_{n}]\,,
%\end{equation*}
%then let 
%$$\beta\defeq[0;b_1,\ldots,b_n]\,, \quad \textrm{and}\quad \gamma\defeq[0;c_1,\ldots,c_{n}]\,,$$ 
%and stop.
\item If the algorithm does not stop after finitely many steps, let 
$$\beta\defeq[0;b_1,b_2,\ldots]\,, \quad \textrm{and}\quad \gamma\defeq[0;c_1,c_2,\ldots]\,.$$ 
\end{enumerate}
\end{defn}

\begin{prop}
For any field $\fq$ consider the infinite continued fractions
\begin{equation}
\label{beta_gamma_def}
\beta=[0;t,t^5,t^9,t^{13},\ldots]\quad\text{and}\quad
\gamma=[0;t^3,t^7,t^{11},t^{15},\ldots]\,,
\end{equation}
and put $\alpha=\beta+\gamma$. Then the decomposition of $\alpha$ as in Definition \ref{def:shulga_an} gives back $\beta$ and $\gamma$.    
\end{prop}

\begin{proof}
For the purpose of this proof, let $b_n=t^{-(4n-3)}$ and $c_n=t^{-(4n-1)}$ for every $n\geq1$, and let $\frac{p_n}{q_n}$ and $\frac{s_n}{t_n}$ be the $n$th convergents of $\beta$ and $\gamma$, respectively, for every $n\geq0$. %We will show that the numbers $\beta$ and $\gamma$ in the Definition \ref{def:shulga_an} coincide with the numbers defined in \eqref{beta_gamma_def}. 
 To verify the proposition, it is sufficient to show that
\begin{equation}
\begin{split}
\label{bn_cn_example}
b_j=&a_j\left(\alpha-\frac{s_{j-1}}{t_{j-1}}\right),\\
c_j=&a_j\left(\alpha-\frac{p_{j}}{q_{j}}\right)
\end{split}
\end{equation}
for all $j\ge 1$.
%, where $b_n$ and $c_n$ are as in Definition \ref{def:shulga_an}. 
We verify \eqref{bn_cn_example} by complete induction on $j$. 
%Suppose that $n=1$, then by the definition $\frac{s_{n-1}}{t_{n-1}}=0$. The first equality in \eqref{bn_cn_example} is equivalent to the fact that
%\begin{equation}
%\label{b1_def}
%a_1(\beta+\gamma)=t\,.    
%\end{equation}
%Taking into account \eqref{eq:common4}, \eqref{b1_def} is equivalent to
%\begin{equation}
%\label{ex_step1_prop_b1}
%\left|\beta+\gamma-\frac{1}{t}\right|<\frac{1}{|t|^2}\,.
%\end{equation}
%The inequality \eqref{ex_step1_prop_b1} is true due to the ultrametric property of the norm in $\fttt$. Indeed,
%\begin{equation}
%\label{ex_step1_prop_ultrametric_b1}
%\left|\beta+\gamma-\frac{1}{t}\right|\le
%\max\left\{\left|\beta-\frac{1}{t}\right|,\left|\gamma\right| \right\}=\max\left\{\frac{1}{|t|\,|t^6+1|}, \frac{1}{|t^3|}\right\}<\frac{1}{|t|^2}\,.
%\end{equation}
%In the general case, the argument is similar. 
Suppose $n\geq1$ and \eqref{bn_cn_example} holds for every $1\leq j<n$. Due to \eqref{eq:common3} and \eqref{eq:common4}, the first equality in \eqref{bn_cn_example} for $j=n$ is equivalent to 
\begin{equation}
\label{ex_step1_prop}
\left|\alpha-\frac{s_{n-1}}{t_{n-1}}-\frac{p_n}{q_n}\right|<\frac{1}{|q_n|^2}\,.
\end{equation}
Using the ultrametric property of the norm in $\fttt$, we see that
\begin{equation}
\label{ex_step1_prop_ultrametric}
\left|\alpha-\frac{s_{n-1}}{t_{n-1}}-\frac{p_n}{q_n}\right|\le
\max\left\{\left|\beta-\frac{p_n}{q_n}\right|,\left|\gamma-\frac{s_{n-1}}{t_{n-1}}\right| \right\}=\max\left\{\frac{1}{|q_n|\,|q_{n+1}|}, \frac{1}{|t_{n-1}|\,|t_n|}\right\}.
\end{equation}
Comparing \eqref{ex_step1_prop} and \eqref{ex_step1_prop_ultrametric}, we see that it is sufficient to verify the inequality
\begin{equation}
\label{2q_n_verify}
\deg(q_n) - \deg(t_{n-1}) < \deg(t_n) - \deg(q_n) \,.
\end{equation}
One can easily see that for every $n\ge 1$ the left side of \eqref{2q_n_verify} is $2n-1$ and the right side is $2n$.

%Now suppose that both inequalities in \eqref{bn_cn_example} hold for $1,2,\ldots,n-1$ and the first inequality holds for the index $n$. 
The second equality in \eqref{bn_cn_example} for $j=n$ is 
%now established using a similar argument: it is 
equivalent to
\begin{equation}
\label{ex_step2_prop}
\left|\alpha-\frac{p_{n}}{q_{n}}-\frac{s_n}{t_n}\right|<\frac{1}{|t_n|^2}\,.
\end{equation}
Applying the ultrametric property once again, we have
\begin{equation}
\label{ex_step2_prop_ultrametric}
\left|\alpha-\frac{p_{n}}{q_{n}}-\frac{s_n}{t_n}\right|\le
\max\left\{\left|\beta-\frac{p_n}{q_n}\right|,\left|\gamma-\frac{s_{n}}{t_{n}}\right| \right\}=\max\left\{\frac{1}{|q_n|\,|q_{n+1}|}, \frac{1}{|t_{n}|\,|t_{n+1}|}\right\}.
\end{equation}
Comparing \eqref{ex_step2_prop} and \eqref{ex_step2_prop_ultrametric}, we see that it is sufficient to verify the inequality
\begin{equation}
\label{2t_n_verify}
\deg(t_n) - \deg(q_n) < \deg(q_{n+1}) - \deg(t_n)\,.   
\end{equation}
One can easily see that for every $n\ge 1$ the left side of  \eqref{2t_n_verify} is $2n$ and the right side is $2n+1$.
\end{proof}

\end{document}
